% MAIN SNAPSHOT — LSPV/RMV::TWOTREEBRIDGE1::2026-10-14T07:01:00+09:00
% !TEX program = pdflatex
\documentclass[11pt]{article}
\usepackage[T1]{fontenc}
\usepackage{lmodern,amsmath,amssymb,amsthm,booktabs,tabularx,xurl,hyperref}
\usepackage[margin=0.9in]{geometry}
\hypersetup{hidelinks}
\setlength{\emergencystretch}{3em}
\newtheorem{theorem}{Theorem}
\newtheorem{lemma}[theorem]{Lemma}
\newtheorem{corollary}[theorem]{Corollary}
\newcommand{\Col}{\{0,1,2\}}
\title{TWOTREEBRIDGE1: Shared-Circuit Elimination and an Exact Two-Tree Compatibility Bridge}
\author{UNIVERSAL\_CORE\_SYNTHESIS / LSPV-5}
\date{14 October 2026, 07:01 KST}
\begin{document}\maketitle
\begin{center}\fbox{\parbox{0.94\linewidth}{
\textbf{P versus NP remains OPEN.} CORRELATIONELIM1 is not a complete
separation proof. Its strongest new result is an easy-family counterexample
to a proposed representation-to-time implication. This run closes its next
local construction: existentially eliminating one shared three-valued variable
from arbitrary Boolean circuit summaries requires at most three cofactors,
with an exact shared-DAG and witness ledger. It also proves a stronger
source-level interface: proper 3-colouring reduces to compatibility of just
two explicitly given trees, each individually admitting a linear-size summary
circuit. The new project bridge is \textbf{CANDIDATE}, complete proof supplied,
independent review pending. NP-completeness is not a superpolynomial lower bound.}}
\end{center}

\section{The missing implication, without changing the claim}
The desired conclusion is that one fixed NP-complete language $L$ is outside
P. For correct total deterministic uniform deciders, its resource form is
\[
 \forall A\text{ deciding }L\;\forall k\in\mathbb N\;\forall C>0\;
 \forall N_0\;\exists x:\quad |x|\ge N_0,\quad
 T_A(x)>C(|x|+2)^k.
\]
No such claim has been established. Finite correctness tests, exponential
tables, essential factor arity and named proof-system lower bounds do not
supply this quantifier. The hard inputs here may depend on $A$. This differs
from a common-input obstruction for every eligible short program in UWI.

The active source is ROOTGLUE1's fixed total graph encoding and its exact
$\mathbb F_4$ three-root compatibility system. Its tree carrier and affine
elimination remain unchanged. Labels $0,1,2$ in this report name the three
\emph{nonzero} field colors; they do not permit zero as a field color.
The new two-tree encoding is specified below, including rejection outside the
domain. The original source, intermediate circuits, chosen colors and lifted
witnesses are the same realization. Arbitrary polynomial workspace, full
source access, recomputation, compact algebra and alternate algorithms remain
permitted. No universal implementation compiler is imposed on this direct route.

\section{Primary literature and exact transfer scope}
Darwiche and Marquis distinguish conditioning, singleton forgetting, general
forgetting and the representations that support them \cite{dm}. Their
Table 7 on printed page 242 was visually inspected: singleton forgetting for
general NNF is different from efficient elimination of an unbounded set of
variables. The construction below is elementary Shannon expansion, not a
new foundational circuit algorithm. Graph 3-colourability is NP-complete;
Gasarch provides a primary reduction paper and attributes the classical
result to Stockmeyer and Lovasz \cite{gasarch}. The reduction to two trees is
given in full here, rather than inferred from a similarly named graph problem.
De Colnet and Mengel's result concerns DNNF size and regular Resolution,
not arbitrary computation \cite{colnet}. No ETH, GRH, cryptographic or
random-oracle assumption is used.

\section{A complete local shared-DAG elimination theorem}
Let $H(v,z)$ be a fan-in-two Boolean circuit with $s$ Boolean operation gates;
$v$ ranges over $\Col$. The three input bits for $v$ are the equality tests
$[v=0],[v=1],[v=2]$. Other color inputs have the same encoding. Let $s_1$
be the number of gates reachable from the output whose syntactic dependency
set includes $v$, and put $s_0=s-s_1$.

\begin{theorem}[Charged cofactor construction]\label{thm:cofactor}
One deterministic uniform construction returns circuits
$H_0(z),H_1(z),H_2(z)$ and
\[
 G(z)=H_0(z)\lor H_1(z)\lor H_2(z)
      \equiv\exists v\in\Col\ H(v,z).
\]
The shared DAG for these four outputs has at most $s_0+3s_1+2$ gates.
In an append-only pool that already stores $H$, at most $3s_1+2$ new gates
are allocated. Evaluation of the three cofactor outputs returns all valid
choices of $v$. A supplied satisfying assignment to $G$ therefore lifts to a
satisfying assignment to $H$ in polynomial time in the actual circuit length.
\end{theorem}
\begin{proof}
Compute dependency sets by one topological pass, charging their bit storage.
Reuse every gate not depending on $v$. For each $c\in\Col$, replace all
three equality input bits of $v$ by the corresponding constants, and copy
only the $v$-dependent gates in topological order. Each copy uses the same
cofactor in both children. The resulting output is $H_c=H|_{v=c}$ by
induction over gates. Two OR gates combine the three outputs. This proves
the size and allocation bounds; constant propagation or syntactic interning
can only reduce them. Indices and dependency sets have polynomial bit length;
a deterministic scan implementation has polynomial construction time and space,
without assuming unit-cost arbitrary-width words or free hashing.

For any fixed $z$, evaluate the three outputs and return any $c$ with
$H_c(z)=1$. Such a color exists exactly when $G(z)=1$. In particular, for
$H=H_1^{\rm source}\land H_2^{\rm source}\land H_3^{\rm source}$, every
source summary uses the \emph{same} selected color, not three independent
existential choices. Internal witness lifts for the summaries compose when
all shared variables are fixed and their remaining private variables are
disjoint. Otherwise those variables must be exposed as part of the interface.
\end{proof}

This closes the prior one-variable/three-summary assignment at greater
generality than trees. It does not prove efficient elimination of all
variables. If $S_j$ is active size, the unconditional construction only gives
\[
 S_{j+1}\le3S_j+2,\qquad S_j\le3^j(S_0+1)-1.
\]
These are \emph{upper} bounds, not lower bounds. The physical archive of
cofactor outputs used in reverse reconstruction is charged separately; a
simple bound is $O(3^j(S_0+1))$ allocated gates. An arbitrary solver is not
required to retain this archive or perform these expansions. In particular,
for a satisfiable instance the final scalar answer is true, irrespective of
the size of this specific construction history.

\section{Two explicit forests already carry arbitrary 3-colouring}
Let $\Gamma=([n],E)$ be a simple graph, $n\ge1$, with $m$ edges. For each
vertex $v$, make $d'_v=\max(1,\deg v)$ occurrence vertices, one for each
incident edge, or one for an isolated vertex. Set
\[
 s=\sum_v d'_v=2m+n_0,\qquad t=s-n,
\]
where $n_0$ counts isolated source vertices. Consecutive occurrences $x,y$
of the same source vertex are connected by fresh vertices $a,b$ and the
five-edge equality diamond
\[
 xa,\ xb,\ ab,\ ay,\ by.
\]
In a proper 3-colouring, $a,b$ differ, so both $x,y$ must take the remaining
color. Conversely, if $x=y$, assign the other two colors to $a,b$ in either
order. Hence this gadget is an exact equality relation, with a polynomial
source-witness extension, not a unique extension.

Partition its edges as follows:
\[
 E_1^{\rm gadget}=\{xa,ab,by\},\qquad
 E_2^{\rm gadget}=\{xb,ay\}.
\]
Put the edge corresponding to each original source edge between its two
occurrences into $E_2$. There are $q=s+2t=3s-2n$ vertices in the resulting
shared graph $J$.

\begin{lemma}[The partition really consists of forests]\label{lem:forests}
$(V(J),E_1)$ is a disjoint union of $n$ paths. $(V(J),E_2)$ is a forest
with $s-m$ components. Every proper 3-colouring of $J$ projects to a proper
3-colouring of $\Gamma$, and every source coloring extends to $J$.
\end{lemma}
\begin{proof}
The three edges in $E_1^{\rm gadget}$ are the path $x-a-b-y$;
concatenating equality gadgets makes one path per source vertex, including
isolated paths. In $E_2$ each $a$ or $b$ is a leaf attached to an occurrence.
After removing those leaves, all original-edge links form a matching,
since each occurrence was reserved for exactly one source edge. There can
be no cycle. Counting edges gives $|E_1|=3t$ and $|E_2|=2t+m$, so the
component counts are $q-3t=n$ and $q-(2t+m)=s-m$.
The diamond equality implies all occurrences of a source vertex agree.
Original-edge links then enforce exactly the source edge inequalities.
Conversely the explicit other-two-colors assignment extends every source
coloring. These operations are polynomial in $n+m$.
\end{proof}

\section{Connecting each forest without adding a constraint}
For each forest independently add a private hub $h_i$. For every component
choose a representative $r$, add a private vertex $p_r$, and add the path
$r-p_r-h_i$. This connects the components without cycles. Regardless of the
color of $r$ and $h_i$, one of three colors differs from both, so every
coloring of the forest extends. Fixing the hub color, for example to label
zero during the witness construction, is legal and takes no advice.

The resulting trees $T_1,T_2$ share exactly the original $q$ vertices of
$J$. All added hubs and connectors are private. If $k_1=n$ and $k_2=s-m$,
then
\[
 |V(T_i)|=q+k_i+1,\qquad |E(T_i)|=q+k_i,
\]
and the size of their union is
\[
 h=q+k_1+k_2+2=4s-n-m+2=O(n+m).
\]
Shared vertices are arbitrary vertices of the trees, \emph{not necessarily
leaves}. The prior all-leaf projection-count theorem is not being applied
outside its hypothesis. The tree dynamic-programming compiler itself permits
any designated boundary subset.

\begin{theorem}[Two-tree compatibility bridge]\label{thm:twotree}
There is a deterministic polynomial-time reduction from graph 3-colourability
to compatibility of two explicitly given trees. A source coloring extends
in polynomial time to colorings of both trees agreeing on their shared
vertices. Any such pair projects in polynomial time to a source coloring.
Each tree has an individually constructible $O(|V(T_i)|)$-gate Boolean
summary $R_i(z)$ on the shared boundary. Nevertheless deciding
\[
 \exists z\in\Col^{V(T_1)\cap V(T_2)}
             [R_1(z)\land R_2(z)]
\]
is NP-complete on the stated total encoding.
\end{theorem}
\begin{proof}
Apply Lemma~\ref{lem:forests} and the private hub construction. Projection
restricts to the shared graph and takes one occurrence per source vertex.
Extension assigns the diamond vertices, then the private hubs and connectors.
All original constraints are checked on this same witness. These maps preserve
satisfiability in both directions but need not preserve the number of witnesses.

For a tree summary, root the tree and define for each color $c$
\[
 A_v(c)=I_v(c)\land\bigwedge_{u\text{ child of }v}
                           \bigvee_{a\ne c}A_u(a).
\]
For a shared vertex $v$, $I_v(c)$ is its boundary equality input; for a private
vertex it is true. The root disjunction is the exact extension predicate,
with a constant number of gates per edge. The tree's explicit decomposition
and the compiler are provided, not acquired by solving an unknown recognition
problem. Full boundary assignments can also be extended by reverse traversal.
A coloring of the finite union is a polynomially checkable NP certificate.
The preceding reduction and known 3-colouring completeness give NP-hardness.
\end{proof}

\paragraph{A total, fixed language.}
Use header \texttt{01010100}, unary universe size $1^h0$ with $h\ge1$, and,
for each of the two trees, $h$ membership bits followed by $\binom h2$
upper-triangular adjacency bits. Require canonical lengths, no edges outside
the declared vertex set, each graph a nonempty connected tree, and every
universe label used by at least one tree. Reject every other string.
This length is exactly
\[
 N=9+h^2+2h.
\]
A certificate consists of $2h$ color bits with the unused fourth code rejected.
Both tree constraints can be checked in polynomial time. The reduction has
polynomial output length; invalid source encodings can be mapped to a fixed
invalid target encoding. This is a fixed NP-complete language with YES and NO
inputs, not one formula, a promise-only hidden family, or an all-NO benchmark.

\paragraph{What this theorem does not say.}
The conjunction $R_1\land R_2$ itself has a small circuit, simply their two
circuits plus an AND gate. The difficult target is \emph{finding a satisfying
boundary assignment or proving none exists}, not constructing that conjunction
or evaluating it on a supplied assignment. NP-completeness does not prove
superpolynomial time. This bridge relates the source to actual overlapping
summaries, but does not resolve the original all-algorithm lower bound.

\section{Why the proposed contradiction under P=NP does not close}
A useful hypothesis audit gives a stronger repair than multiplying local
circuit growth bounds. Let $I$ be the original fixed-language input and let
$\mu$ assign each vertex either a color or ``unassigned.'' Define
\[
 \operatorname{EXT}(I,\mu)=1
 \iff\text{a valid global witness extending }\mu\text{ exists}.
\]
This is an NP predicate with polynomial witness length and a total
polynomial verifier. Thus, \emph{under the assumption P=NP}, some fixed
uniform polynomial-time program $D$ decides it.

For each input length, unroll $D$'s bounded computation into an ordinary
Boolean circuit. The straightforward space-time tableau construction has
polynomial size and polynomial construction cost; for a $T$-step machine a
conservative quadratic-size simulation suffices. Hardwire $I$ and retain the
pin bits $\mu$ as circuit inputs. One polynomial-size circuit then answers
\emph{all} partial-assignment extension queries for that same $I$. Every
elimination prefix is a special case of this single predicate. Witness search
uses at most three queries per vertex and verifies the final answer.

This is a conditional consequence of P=NP, not a new unconditional algorithm.
It neither restricts the circuit to DNNF nor guarantees a small-space or
near-linear runtime. Therefore it contradicts none of the retained
representation, proof-system or joint time-space bounds. In particular,
repeatedly expanding one chosen cofactor representation cannot disprove the
hypothesis while this alternative is left unaddressed. Calling an assertion
that all such alternative computations fail a ``compression lemma'' would
merely rename the missing separation-strength premise.

A direct lower bound for the fixed language need not follow this circuit
route. The present audit does not demand a universal implementation compiler
or a general circuit lower bound as a mandatory approach.

\section{Executed audit and exact limitations}
The standard-library script
\path{experiments/review/two_tree_bridge_audit.py}
has SHA-256
\path{70d2e28d14ed371ce4d0783ba629ace92b95a6a1aabdee39b7841530f9bf3554}.
Fourteen separately bounded invocations passed \textbf{910,185 assertions},
zero failures. The maximum invocation took 12.089 seconds rounded upward;
the sum was 36.811 seconds rounded upward. Each invocation has a 25-second
audit deadline. The four degree-five reduction chunks each contain 256 graphs.
No production routine has a fixed input-size cap.

\begin{center}\small
\begin{tabularx}{\linewidth}{p{0.28\linewidth}X}
\toprule
Check & Executed scope\\\midrule
Circuit elimination & All 1,099 labelled graphs on 1--5 vertices;
126,577 exact prefix/assignment projection comparisons; 5,405 construction
and retained-support gate bounds each; 1,032 returned witnesses.\\
Two-tree reduction & The same 1,099 source graphs; 254,253 original color
assignments with canonical gadget extensions and exact compiled-summary checks;
2,198 actual tree validations and 1,099 exact encoding checks.\\
Independent target search & DSATUR backtracking on 77 constructed union graphs,
without access to source/gadget semantics; 75 positive witnesses projected
back to the source, two negative decisions.\\
Small controls & All diamond endpoint choices, all hub endpoint choices,
one variable genuinely shared by three tree summaries, internal witness lifts,
and six malformed target records.\\
\bottomrule\end{tabularx}\end{center}

The largest reduction has 67 union vertices and 50 shared variables; the
largest individual summary uses 329 gates and the largest conjunction 644.
Those small finite sizes do not extrapolate to a polynomial global solver.
The target tests exhaust all original source assignments, not all assignments
to the much larger target graph. The correctness of every target-to-source
projection follows from the displayed diamond proof; independent target search
provides an additional finite check, not an exhaustive asymptotic argument.
The circuit tests include both YES and NO graphs and compare to complete
source model enumeration, not to a table generated by the same cofactor code.
All finite results are \textbf{COMPUTATIONAL}. New project theorems remain
\textbf{CANDIDATE} pending independent review.

\section{Corpus, preservation and publication discipline}
\paragraph{CORPUS REVIEWED.}
Read the current root, current AGENTS write/recovery rules, latest specialist
commit history and OUTQUAD64 sentinel. Used the exact uploaded ROOTGLUE1 and
CORRELATIONELIM1 sources, code, raw data and root bytes, not reconstructed
scientific text. The full research skill, CODEX, seven task specifications,
main4/main5 drivers and latest four actual specialist sources were already
read in this conversation and their recorded scopes were retained. A complete
new reread of all historical includes or original Drive passages is not claimed.
Direct container GitHub access again failed DNS, so repository-wide checkout
preflight/CI was not executed. Targeted code, proof, artifact and readback
checks are stated separately.

\paragraph{DUPLICATION / CONTRADICTION AUDIT.}
Repository searches for ``two forests'' returned no matching source in the
bounded search. The ``cofactor'' search recovered LOGCOF1 and its actual
handoff: logarithmically many withheld bits permit polynomial precomputation
with all work charged. This is consistent with the present cofactor accounting,
not a forgotten proof of superpolynomial work. General Shannon expansion and
tree compilation are known. No broad novelty priority is claimed for the
elementary two-forest reduction. Its value here is its exact, tested interface
with the project's compact summary construction.

\paragraph{WHAT CHANGED RELATIVE TO ROOT MAIN.}
The one-variable/three-summary target is completed by a general cofactor
construction with charged physical sharing and support retention. The next
source target is reduced to two tree summaries, not arbitrarily many unknown
algebraic objects. The joint language is explicitly total and NP-complete.
The hypothesis P=NP is audited against one source-fixed extension circuit,
showing precisely why the apparent expansion contradiction is not complete.
No unrestricted lower-bound status is promoted.

\paragraph{ADOPTED / REJECTED / RETIRED PRIOR COMPONENTS.}
Adopted: affine elimination, root-domain maps, exact projected relations,
input-assisted witness lifting, unrestricted source access and polynomial
workspace. Rejected: inferring a lower bound from an upper-bound recurrence,
identifying small-circuit evaluation with finding a satisfying assignment,
and claiming a single bounded forgetting step has unavoidable exponential
cost. Preserved: all valid historical selector algorithms and restricted-model
results at their hypotheses, plus the conditional UWI amplification implication.

\paragraph{EARLY22/SINGLE22 acknowledgment.}
Already integrated through the historical ledger, not newly reproved. Retain
separately the signed exact-$(2,2)$ compiled minimal-UNSAT/Resolution family,
the clause-monotone cubic-pair theorem/calibrations, and affine continuation.
LBO3SAT-0031/0032 and PROOF-0034 retain their assumptions; historical $N$
counts variables, not encoded bits. This new two-tree source does not inherit
PC22 minimality, recognition, or proof projection. Its independent review
is pending. A direct non-P theorem for either correctly defined NP-complete
language needs no extra broadening step.

\paragraph{Actual main4/main5 and predecessor status.}
main4 remains the achievable-region driver; main5 the WRAPPER construction
driver. Their conditional acquisition/replay ledger and open all-input coverage
are not changed. OUTQUAD64 is the latest located substantive Bridge report,
with header/sentinel and prior integration intact. It stays candidate, and
its 64-candidate hidden-map theorem is not used to solve this joint language.
No exact missing body is available or reconstructed. The root restored during
ROOTGLUE1 remains substantive; no duplicate recovery report is created.
Existing schedules and team membership are unchanged. Fresh specialist
acknowledgments are requested, not invented.

\section{Barrier audit and the surviving proof obligation}
The new reduction is a constructive many-one transfer on explicit input,
not an inference from a geometric or cryptographic analogy. Its witnesses,
bit encodings and all maps have polynomial size. The compact circuit
construction charges the full shared pool and retained cofactor outputs.
Auxiliary variables, other circuit encodings, reordered operations, memoization
and recomputation remain legal alternatives for a decider. Program length
is not its workspace. DNNF/PC/PCR/Resolution lower bounds are still not
whole-computation results. The cofactor routine and easy-tree routines survive
stronger semantic and algebraic oracle access; no nonrelativizing or
nonalgebrizing contradiction has been obtained. Natural Proofs is not invoked.
Post-quantum input-output relation spaces remain useful algorithmic attacks,
not unconditional hardness premises; OUTQUAD64 is preserved on that basis.

The shortest surviving chain is
\[
 \mathrm{3COL}\ \le_m^{\rm P}\ \mathrm{TWO\mbox{-}TREE\mbox{-}COMPATIBILITY}
 \quad\text{with exact polynomial witness maps},
\]
followed by the still \textbf{OPEN} arrow to unrestricted superpolynomial
runtime. The latter is not a minor technicality. Alternatively, the retained
UWI benchmark
\[
 H(F)=\min_{|p|\le b_0+\lfloor\log_2(|F|+2)\rfloor}
                   (t_p(F)+q(|F|)+1)
\]
would be superpolynomial if a satisfiable sequence had
$s_j\to\infty$, $s_{j+1}\le As_j$ and
$H(F_{j+1})\ge H(F_j)r(s_j)/B$ with fixed $A,B$ and $r(s)\to\infty$.
This implication is conditional; neither the two-tree bridge nor these
experiments establish its premise. The common-input quantifier remains distinct.

\textbf{One next assignment.} Audit repeated shared-DAG elimination on the
explicit two-tree reductions, with all cofactor supports and source-witness
lifts charged, and compare source-aware recomputation before proposing a
resource invariant. Bridge reviews the exact two-tree and off-domain maps;
Tests independently checks target witnesses; Redteam attacks quantifiers and
small-circuit-versus-search claims; Lower Bound must give an independently
justified work measure before asserting unavoidable growth. No separation
claim is authorized by a failed chosen compiler.

\begin{thebibliography}{9}
\bibitem{dm} A. Darwiche and P. Marquis. A Knowledge Compilation Map.
Journal of Artificial Intelligence Research 17, 229--264, 2002.
\url{https://arxiv.org/abs/1106.1819}.
\bibitem{gasarch} W. Gasarch. A Sane Proof that $COL_k\le COL_3$.
arXiv:1407.5128v2, 2016. \url{https://arxiv.org/abs/1407.5128}.
\bibitem{colnet} A. de Colnet and S. Mengel. Characterizing Tseitin-Formulas
with Short Regular Resolution Refutations. JAIR 76, 265--286, 2023.
\url{https://www.jair.org/index.php/jair/article/view/13521}.
\end{thebibliography}
\end{document}
% END_OF_VERIFIED_MAIN::TWOTREEBRIDGE1
